\chapter{Experimental Setup and Simulation Environment}
\label{ch:experimental}

Chapter~\ref{ch:methodology} presented the architecture of the four-stage pipeline
and the V-model discipline under which it was built. This chapter specifies the
conditions under which that pipeline is exercised and measured: the datasets that
feed it, the simulation environment that supplies simultaneous multi-agent ground
truth, the concrete model and algorithm configurations of each stage, and the
evaluation metrics against which every stage is scored. The results obtained under
this setup are reported and discussed in Chapter~5; the present chapter establishes
\emph{what} is measured and \emph{how}, so that those results are reproducible and
interpretable.

The experimental design is a direct consequence of the dual-dataset strategy of
Section~\ref{sec:data}. The single-vehicle perception chain (Stages~1--3) is
exercised on real driving imagery from the KITTI benchmark suite, the only source
that can confirm camera-only perception on genuine sensor data. The cooperative
layer (Stage~4) is exercised on the CARLA simulator, the source that can provide
two vehicles observing one scene at one instant with exact relative pose and
occlusion-truthful per-agent visibility. Each dataset is described in the role it
serves, followed by the per-stage configurations and the metric definitions.

\section{Implementation and Software Environment}
\label{sec:environment}

The pipeline is implemented in Python, with each stage as a standalone, command-line
runnable module that communicates with its neighbours only through on-disk artifacts
(disparity maps, detection files, lifted-position files), consistent with the
modularity requirement (NFR1). The principal third-party components are OpenCV
(classical stereo matching and image I/O), PyTorch and the HuggingFace Transformers
library (the learned depth and detection models), NumPy and Shapely (geometry), and
MLflow (experiment tracking).

Two cross-cutting properties of the environment are essential to the validity of the
reported results:

\begin{itemize}[leftmargin=*,itemsep=3pt]
  \item \textbf{Determinism.} The random seeds of the \texttt{random}, NumPy, and
  PyTorch libraries are fixed to $42$ at every entry point. Combined with the
  configuration-driven design --- every operational value (paths, thresholds, model
  choices, tuned parameters) lives in a per-stage YAML configuration file rather
  than in code --- this makes each run reproducible from its recorded configuration
  (NFR2).
  \item \textbf{Experiment tracking.} Every run logs its parameters and resulting
  metrics to a persistent MLflow store, so that no result exists only as transient
  console output and any reported figure can be traced back to the run that produced
  it.
\end{itemize}

All experiments in this work were run on \textbf{CPU only}, with no GPU
acceleration. This constraint is itself a meaningful part of the setup: it reflects
the cost-constrained, camera-based deployment scenario motivating the thesis, and it
directly shapes the runtime trade-off between the classical and learned depth
estimators discussed in Section~\ref{sec:stage1cfg}.

\section{Real-Image Stereo Dataset: KITTI}
\label{sec:kitti}

The single-vehicle chain is developed and validated on the KITTI vision benchmark
suite~\cite{geiger2012kitti}, a standard real-world autonomous-driving dataset
recorded from a stereo camera rig mounted on a moving vehicle. Two of its splits are
used, each matched to the stages it can validate.

\subsection{Object split (Stages 1--2)}

The object detection split provides individual stereo frames with rectified
left/right images, per-frame camera calibration, dense disparity ground truth, and
labelled 2D object boxes. It is used to validate Stage~1 (depth) and Stage~2 (2D
detection) in isolation. The validation reported in Chapter~5 uses the first ten
frames of the split (\texttt{000000}--\texttt{000009}); the disparity ground truth
of this split is what makes a direct, per-pixel depth-accuracy measurement possible.

\subsection{Tracking split (Stage 3)}

Validating the lifting stage requires \emph{3D} ground truth --- the true metric
position of each vehicle --- which the object split does not align with its stereo
frames. Stage~3 is therefore validated on the KITTI tracking split, whose sequences
carry per-frame 3D object labels registered to the camera frame. The Stage~3
evaluation in Chapter~5 uses five sequences (24 frames in total). Each 2D detection
is lifted to a 3D position and compared against the tracking labels by 3D centre
distance (Section~\ref{sec:metricslift}).

\begin{table}[htbp]
\centering
\caption{KITTI conventions adopted by the pipeline.}
\label{tab:kitti}
\renewcommand{\arraystretch}{1.3}
\begin{tabularx}{\textwidth}{@{}lL@{}}
\toprule
\textbf{Item} & \textbf{Convention} \\
\midrule
Calibration keys & \texttt{P2}, \texttt{P3} (left/right projection matrices), \texttt{R\_rect\_00}, \texttt{Tr\_velo\_to\_cam} \\
Stereo baseline  & recovered as $B = (t_x^{L} - t_x^{R})/f$ from \texttt{P2}/\texttt{P3} \\
Disparity ground truth & 16-bit PNG; metric disparity $= \text{raw}/256.0$; $\text{raw}=0$ marks an invalid (no-measurement) pixel \\
Object vertical reference & label $y$ is the \emph{bottom} of the object (centre $= y - h/2$) \\
Coordinate frame & camera optical frame: $X$ right, $Y$ down, $Z$ forward \\
\bottomrule
\end{tabularx}
\end{table}

\subsection{KITTI conventions}

The pipeline adopts the standard KITTI conventions summarised in
Table~\ref{tab:kitti}. Of particular note for the depth and lifting stages: the
disparity ground truth is stored as a 16-bit image scaled by $256$, with a raw value
of zero denoting ``no measurement'' rather than zero disparity --- the explicit
invalid-pixel contract that Stage~1 preserves so that downstream stages never
mistake an absent measurement for a near-infinite depth.

\section{Simultaneous Multi-Agent Dataset: CARLA}
\label{sec:carla}

The cooperative layer cannot be validated on KITTI, which records only one vehicle's
view at a time. Stage~4 is therefore exercised on data generated with the CARLA
simulator~\cite{dosovitskiy2017carla}, which can render two vehicles observing the
same scene at the same instant with exact, noise-free relative pose --- and,
crucially, with per-object, per-agent visibility derived from the rendered view
rather than guessed geometrically.

\subsection{Scenario}

The scenario is a single urban intersection (the CARLA \texttt{Town10HD} map)
recorded over 300 synchronised frames, populated with traffic and two moving ego
vehicles. The two egos form a \textbf{leader--follower} pair: Vehicle~A travels close
to the shared traffic (median distance to the co-observed cars $\approx$~11~m) while
Vehicle~B follows roughly $3.3\times$ further behind (median $\approx$~36~m). This
geometry is deliberate --- it places the same vehicles at very different ranges from
the two agents, which is exactly the heterogeneous-uncertainty condition the thesis
sets out to study, and it furnishes a population of cars that the follower cannot
perceive alone, i.e.\ the blind-spot recovery that cooperation is meant to address.

\subsection{Per-agent recordings and data layout}

Each of the two agents (\texttt{vehicle\_a}, \texttt{vehicle\_b}) contributes, per
frame, a rectified stereo pair (\texttt{image\_left}/\texttt{image\_right}), a camera
calibration (\texttt{calib.json}), a noise-free disparity reference
(\texttt{disp\_noc\_0}), and its world pose (\texttt{pose.json}). A separate
per-frame ground-truth record (\texttt{gt\_boxes}) lists every vehicle in the scene
with its 3D box and the visibility metadata described below. This layout lets the
same Stage~1--3 chain that runs on KITTI run unchanged on each CARLA agent (the
detector path), after which Stage~4 fuses the two agents' outputs using the relative
pose derived from the two \texttt{pose.json} files.

\subsection{Visibility-truthful per-agent and cooperative ground truth}
\label{sec:coopgt}

The defining property of this dataset is that an object's visibility to each agent is
\emph{occlusion-truthful}. Each ground-truth box carries the number of its pixels
actually rendered to each agent's camera; a car counts as seen by an agent only if at
least \texttt{min\_visible\_pixels} ($=10$) of its pixels are visible to that agent.
Cars hidden behind buildings or other vehicles therefore correctly drop out of an
agent's truth set --- a geometric field-of-view test would have falsely credited an
agent with seeing them and thereby understated the cooperative benefit.

From these per-agent truth sets the \textbf{cooperative ground truth} is formed as
the union of every (non-ego) vehicle visible to A \emph{or} B, de-duplicated by
object identity (\texttt{actor\_id}). The two ego vehicles are part of the scene but
are not perception targets --- their poses are shared directly over the cooperative
link --- so they are excluded from the cooperative ground truth, and an agent's
detection of the \emph{other} ego is treated as an \textbf{ignore region} (neither a
true nor a false positive). Over the 150 frames evaluated in Chapter~5, the
cooperative ground truth comprises 9 distinct cars across 586 visible instances.

\section{Model and Algorithm Configurations}
\label{sec:configs}

Each stage is governed by a per-stage YAML configuration. The settings that
materially affect the results are summarised here; the full files are part of the
codebase.

\subsection{Stage 1 --- depth estimation}
\label{sec:stage1cfg}

The two interchangeable estimators of Section~\ref{sec:stage1} are configured as
follows. The classical estimator is OpenCV Semi-Global Block
Matching~\cite{hirschmuller2008} with a 192-pixel disparity search range, an
$11\times11$ matching block, smoothness penalties $P_1=8$ / $P_2=32$, a left--right
consistency check (\texttt{disp12\_max\_diff}~$=1$), a uniqueness ratio of $10\%$,
and speckle filtering --- a configuration that yields a \emph{sparse} but
consistency-checked disparity map. The smoothness penalties are kept intentionally
low because KITTI scenes have sharp depth discontinuities at object boundaries that
high penalties would over-smooth. The learned estimator is
WAFT-Stereo~\cite{wang2026}, run as an in-process inference step from a pretrained
checkpoint, producing a \emph{dense} (full-coverage) disparity map. On the CPU-only
environment of Section~\ref{sec:environment} the two differ sharply in cost ---
SGBM runs in the order of a second per image, WAFT in the order of a minute --- a
trade-off that is itself part of the analysis.

\subsection{Stage 2 --- 2D detection}

Detection uses RT-DETR~\cite{lv2024} (the \texttt{rtdetr\_r50vd} checkpoint,
$\approx$~43\,M parameters) applied to the left image, pretrained on the COCO object
vocabulary. The relevant COCO vehicle categories are remapped onto the project's
single target class --- \texttt{car}, \texttt{truck}, and \texttt{bus} all map to
KITTI \texttt{Car} --- and all other categories are discarded. A confidence
threshold of $0.5$ suppresses weak detections. Consistent with the Car-only scoping
of Section~\ref{sec:stage2}, no pedestrian or cyclist classes are emitted.

\subsection{Stage 3 --- lifting}

Lifting applies the per-method depth-sampling strategy motivated in
Section~\ref{sec:stage3}: for the sparse SGBM map, the $20^{\text{th}}$ percentile of
valid depths within a detection region with no vertical crop; for the dense WAFT map,
the $35^{\text{th}}$ percentile over the top $40\%$ of the box with depths nearer than
6~m gated out as road surface. A detection with fewer than 10 valid depth pixels in
its region is skipped rather than lifted from an empty frustum. These values are a
single source of truth in the Stage~3 configuration, read identically by the
pipeline and its validator.

\subsection{Stage 4 --- cooperative fusion}

Fusion matches A's observations against B's (registered into A's frame) greedily by
bird's-eye-view centre distance within class, accepting a pair as the same object
when their separation is at most $\tau_{\text{Car}}=2.0$~m. A matched pair is fused
only if its post-registration BEV displacement is below 1.0~m; a larger displacement
signals a faulty match or pose error (not motion, since the observations are
simultaneous) and the pair is flagged and kept unmerged. Fused confidence follows the
noisy-OR rule and fused position the confidence-weighted average of
Section~\ref{sec:stage4}.

\begin{table}[htbp]
\centering
\caption{Summary of the per-stage experimental configuration.}
\label{tab:configsummary}
\renewcommand{\arraystretch}{1.3}
\begin{tabularx}{\textwidth}{@{}llL@{}}
\toprule
\textbf{Stage} & \textbf{Component} & \textbf{Key settings} \\
\midrule
1 & SGBM / WAFT-Stereo & 192-px range, $11\times11$ block, $P_1/P_2=8/32$, LR-check / pretrained dense model \\
2 & RT-DETR (\texttt{r50vd}) & COCO $\rightarrow$ Car (car/truck/bus), confidence $\ge 0.5$ \\
3 & Lift-to-position & SGBM: p20, no crop; WAFT: p35, top-40\% crop, 6\,m gate; skip $<10$ px \\
4 & Late fusion & BEV match $\tau_{\text{Car}}=2.0$\,m, fuse if displacement $<1.0$\,m, noisy-OR confidence \\
\bottomrule
\end{tabularx}
\end{table}

\section{Evaluation Metrics}
\label{sec:metrics}

Each stage is scored against ground truth in isolation, so that a failure at any
level can be diagnosed without trusting its successors (NFR3). The metrics for each
stage are defined below.

\subsection{Stage 1 --- depth}

Disparity quality is measured against the dense ground-truth disparity, over valid
pixels only, by two standard KITTI metrics and a coverage figure:

\begin{itemize}[leftmargin=*,itemsep=3pt]
  \item \textbf{End-Point Error (EPE)} --- the mean absolute disparity error in
  pixels, $\text{EPE} = \frac{1}{N}\sum_i |d_i - d_i^{\text{GT}}|$.
  \item \textbf{D1 error} --- the percentage of pixels whose disparity error exceeds
  the standard KITTI outlier threshold (3~px and $5\%$ of the true disparity),
  capturing gross mismatches rather than sub-pixel noise.
  \item \textbf{Coverage} --- the percentage of pixels for which the estimator
  reports a valid disparity. This distinguishes the sparse SGBM map from the dense
  WAFT map and is essential context for the two preceding metrics, which are computed
  over reported pixels only.
\end{itemize}

\subsection{Stage 2 --- 2D detection}

Detection is scored by the standard average-precision framework. A predicted box is a
true positive when its Intersection-over-Union (IoU) with a ground-truth box of the
same class exceeds $0.5$. Per-class Average Precision (AP) and the mean over classes
(mAP) are reported, together with per-class precision and recall at the fixed
operating confidence threshold of $0.5$. Average precision is computed by 11-point
interpolation; a reliability flag marks evaluations made on fewer than ten frames as
small-sample estimates.

\subsection{Stage 3 --- lifting to 3D position}
\label{sec:metricslift}

Because Stage~3 emits a 3D \emph{position} plus the source 2D box --- not a full 3D
box --- it cannot be scored by 3D IoU. Instead a predicted position is matched to the
nearest unmatched same-class ground-truth object within a 3D centre-distance
threshold (Car~$\le 2.0$~m), greedily. From the resulting matches the pipeline
reports precision, recall, and their harmonic mean $F_1$, alongside two localisation
errors averaged over matched pairs: the \textbf{depth error} (error along the camera
forward axis $Z$) and the \textbf{3D centre distance}. To expose the central
range-dependence claim of Chapter~\ref{ch:introduction}, errors are additionally
broken down into depth-range bins (0--10, 10--20, 20--40, and beyond~40~m).

\subsection{Stage 4 --- cooperative fusion}
\label{sec:metricsfusion}

Stage~4 is evaluated against the cooperative ground truth of
Section~\ref{sec:coopgt}. For each frame, three prediction sets are scored against
the \emph{same} cooperative ground truth: \textbf{A-alone}, \textbf{B-alone} (B's
single-vehicle output registered into A's frame), and \textbf{Fused}. Each set is
matched to the cooperative ground truth by the same greedy BEV centre-distance rule
($\tau_{\text{Car}}=2.0$~m), giving per-set precision, recall, and a bird's-eye-view
\textbf{localisation error} over true positives. Beyond these per-set figures, the
following metrics characterise the cooperation itself:

\begin{itemize}[leftmargin=*,itemsep=3pt]
  \item \textbf{Symmetric cooperative gain.} Because the fused scene is shared, the
  benefit is measured for both agents against the same ground truth: A's gain from B
  is the set of cooperative-GT objects only B saw (\texttt{b\_unique\_tp}), and B's
  gain from A is the set only A saw (\texttt{a\_unique\_tp}), each reported as a
  recall improvement and a count of recovered objects.
  \item \textbf{Corroboration rate.} Of the objects co-observed by both agents, the
  fraction that are actually merged into a single fused estimate (as opposed to
  surviving as two separate detections). This distinguishes a genuinely
  deduplicated scene from one that is merely a union of the two agents' detections.
  \item \textbf{Angular separation.} The angle subtended at a co-observed car between
  the two agents' viewing directions, reported as a diagnostic for whether
  cross-view triangulation (which would tighten depth) is geometrically available.
  \item \textbf{Precision--recall sweep and localisation-vs-range curve.} A PR sweep
  over the confidence threshold traces the precision/recall trade-off of the three
  sets, and a plot of localisation error against ground-truth range exposes how the
  cooperative localisation behaves as targets recede.
\end{itemize}

\section{Experimental Protocol and Scope}
\label{sec:protocol}

The protocol follows directly from the V-model verification strategy: each stage is
run and scored standalone against the ground truth of the dataset best able to test
it, and only then is the full chain run end-to-end for a cooperating pair. Throughout,
the two depth estimators (SGBM and WAFT-Stereo) are carried through the entire
pipeline in parallel, so that the influence of depth density and accuracy on every
downstream stage --- detection-independent lifting quality and, ultimately,
cooperative recall and localisation --- can be attributed rather than presupposed.

Two scope conditions bound the interpretation of the results in Chapter~5 and are
stated here in the interest of honesty about what the setup can and cannot establish.
First, the cooperative evaluation rests on a \emph{single} scenario --- one
\texttt{Town10HD} intersection with a leader--follower pair and nine Car-only
cooperative-GT vehicles within 40~m --- so it is a characterised feasibility study,
not a multi-scenario benchmark. Second, the leader--follower geometry fixes the two
agents at a narrow angular separation from the shared cars; testing the
triangulation regime, in which two near-perpendicular viewpoints could tighten depth,
would require a different scenario geometry. These conditions do not weaken the
measurement; they delimit precisely the operating regime within which it is valid.

\bigskip
\noindent This chapter has specified the experimental setup: a deterministic, CPU-only,
configuration-driven, MLflow-tracked implementation; the KITTI benchmark suite as the
real-image source for the single-vehicle chain (object split for depth and detection,
tracking split for lifting); the CARLA \texttt{Town10HD} leader--follower scenario as
the simultaneous, occlusion-truthful multi-agent source for cooperative fusion; the
concrete per-stage model and algorithm configurations; and the metric definitions
against which each stage is scored, from per-pixel disparity error through cooperative
recall, localisation, and corroboration. With the setup established, Chapter~5
presents the stagewise validation results and the characterisation of the cooperative
localisation dynamics that this framework was built to measure.
